\section{Research}
\resumeSubHeadingListStart

        \resumeSubheading
        {Temporal Knowledge Leakage in LLM Prediction Tasks}{Northwestern University}
        {Advisor: Prof. Bradly C. Stadie}{Sep 2025 – Present}
        \resumeItemListStart
        \item\small
        Formalized \textbf{temporal knowledge leakage}—the inadvertent use of information not publicly available at a specified reference time—in LLM-based prediction tasks such as legal case outcomes, salary estimation, and stock ranking.
        \item\small
        Introduced a four-phase detection pipeline with well-defined Shapley-weighted metric. 
        \item\small
        Proposed the \textbf{Temporal LLM Agent}, a multi-phase architecture that proactively prevents leakage during generation through iterative claim verification and regeneration.
        \item\small
        Demonstrated substantial leakage reduction while maintaining prediction quality across legal case prediction, salary estimation, and stock ranking benchmarks.
        % \item\small
        % \textit{\textbf{Targeting KDD 2026.}}
        \resumeItemListEnd

        \resumeSubheading
        {LAMP: Linear Attribution Mapping Probe}{Northwestern University}
        {Advisor: Prof. Bradly C. Stadie}{Jan 2024 – Jan 2025}
        \resumeItemListStart
        \item\small
        Developed \textsc{\textbf{L}inear \textbf{A}ttribution \textbf{M}apping \textbf{P}robe (\textbf{LAMP})}, a framework for interpreting black-box language models by fitting locally linear surrogate models grounded in the model's own self-reported explanations.
        \item\small
        Designed a perturbation-based method to probe model prediction sensitivity and extract decision surfaces without accessing gradients, logits, or internal states—enabling practical interpretability for proprietary LLMs.
        \item\small
        Implemented extensive experiments across sentiment classification, controversial-topic detection, and harmful response auditing.
        \item\small
        \textit{\textbf{Accepted as Spotlight at AISTATS 2026.}} \href{https://openreview.net/forum?id=mynHcPXsMv}{(\texttt{Link}: LAMP: Extracting Locally Linear Decision Surfaces from LLM World Models)}
        \resumeItemListEnd

        \resumeSubheading
        {Factual Memorization and Learning in Large Language Models}{Northwestern University}
        {Advisor: Prof. Bradly C. Stadie}{Jan 2024 – Jan 2025}
        \resumeItemListStart
        \item\small
        Conducted a comprehensive analysis of LLM memorization behaviors under different fine-tuning paradigms, including \textsc{SFT} and \textsc{DPO}, building upon and extending insights from \textsc{FinetuneBench}.
        \item\small
        Reproduced key experimental findings from Stanford's \textsc{FinetuneBench}, and systematically evaluated LLM robustness to input rephrasings and temporal shifts in question formulation.
        \item\small
        Proposed a formal distinction between \textbf{passive memorization} (from exposure) and \textbf{positive memorization} (via direct QA supervision), demonstrating that the latter achieves superior memorization efficiency.
        \item\small
        Designed temporal generalization experiments by injecting future-dated facts into the training set. Found no generalization beyond training, but demonstrated that a lightweight system prompt enforcing temporal consistency can mitigate this overfitting.
        % \item\small
        % \textit{\textbf{Qualifying Exam Paper, Northwestern University.}}
        \resumeItemListEnd
        
        \resumeSubheading
        {CUOLR: Unified Off-Policy Learning to Rank}{Princeton University}
        {Mentors: Prof. Mengdi Wang, Prof. Huazheng Wang}{May 2022 – Dec 2023}
        \resumeItemListStart
        \item\small
        Formulated a unified framework that models various click models in off-policy Learning to Rank (LTR) as a \textit{Markov Decision Process (MDP)}, enabling principled application of offline reinforcement learning.
        \item\small
        Proposed the \textsc{\textbf{C}lick model-agnostic \textbf{U}nified \textbf{O}ff-policy \textbf{L}earning to \textbf{R}ank (\textbf{CUOLR})} algorithm that adapts to a wide range of click models \textit{without requiring explicit debiasing or prior knowledge}.
        \item\small
        Demonstrated that offline RL methods (e.g., DQN, SAC, BCQ, CQL) can be applied effectively to LTR by leveraging the MDP formulation, maintaining consistency and robustness under diverse click model assumptions.
        \item\small
        Validated CUOLR on large-scale real-world datasets, achieving state-of-the-art performance and significantly improved robustness across heterogeneous click models.
        \item\small
        \textit{\textbf{Published at NeurIPS 2023.}} \href{https://openreview.net/pdf?id=oDcWnfZyZW}{(\texttt{Link}: Unified Off-Policy Learning to Rank: a Reinforcement Learning Perspective)}
        \resumeItemListEnd
 
\resumeSubHeadingListEnd